%% isometric_layered_device_stack.snippet.tex
%% Archived from: ~/workspace/ResearchOS/[tikz-paper-workflow]/test_quality/isometric_layered_device_stack.tex
%% License: Same as parent manuscript (attributed to user)
%% Purpose: 3D isometric view of stacked dielectric+electrode+trapped-charge device
%% Requires: polymer-paper-preamble.sty (defines iso basis, \IsoBlock, iso callout, iso charge, iso label)
%%
%% Usage:
%%   %% isometric_layered_device_stack.snippet.tex
%% Archived from: ~/workspace/ResearchOS/[tikz-paper-workflow]/test_quality/isometric_layered_device_stack.tex
%% License: Same as parent manuscript (attributed to user)
%% Purpose: 3D isometric view of stacked dielectric+electrode+trapped-charge device
%% Requires: polymer-paper-preamble.sty (defines iso basis, \IsoBlock, iso callout, iso charge, iso label)
%%
%% Usage:
%%   %% isometric_layered_device_stack.snippet.tex
%% Archived from: ~/workspace/ResearchOS/[tikz-paper-workflow]/test_quality/isometric_layered_device_stack.tex
%% License: Same as parent manuscript (attributed to user)
%% Purpose: 3D isometric view of stacked dielectric+electrode+trapped-charge device
%% Requires: polymer-paper-preamble.sty (defines iso basis, \IsoBlock, iso callout, iso charge, iso label)
%%
%% Usage:
%%   %% isometric_layered_device_stack.snippet.tex
%% Archived from: ~/workspace/ResearchOS/[tikz-paper-workflow]/test_quality/isometric_layered_device_stack.tex
%% License: Same as parent manuscript (attributed to user)
%% Purpose: 3D isometric view of stacked dielectric+electrode+trapped-charge device
%% Requires: polymer-paper-preamble.sty (defines iso basis, \IsoBlock, iso callout, iso charge, iso label)
%%
%% Usage:
%%   \input{styles/snippets/isometric_layered_device_stack.snippet.tex}
%% Or copy-paste the TikZ code below into your .tex and edit layer counts, colors, labels, charge positions.

\providecommand{\LayerCountA}{4}  % total layer count
\providecommand{\LayerThicknessA}{0.28}  % bottom electrode thickness (cm)
\providecommand{\LayerThicknessBD}{0.44}  % middle (dielectric) thickness (cm)
\providecommand{\LayerThicknessBCD}{0.34}  % charge layer thickness (cm)
\providecommand{\LayerThicknessBCDE}{0.22}  % top electrode thickness (cm)
\providecommand{\LayerWidthA}{4.8}  % all layers width (cm)
\providecommand{\LayerDepthA}{1.7}  % all layers depth (cm)
\providecommand{\ColorLayerA}{cAmber}  % bottom electrode
\providecommand{\ColorLayerBD}{cTeal}  % dielectric polymer
\providecommand{\ColorLayerBCD}{cBlue}  % charge layer
\providecommand{\ColorLayerBCDE}{cGray}  % top electrode

\begin{tikzpicture}[iso basis]
  \node[font=\fontsize{9}{10.8}\selectfont\bfseries, text=cGray!90!black, anchor=west]
    at (-0.40, -1.10, 2.95) {Layered device stack};

  % Layer A: bottom electrode
  \begin{scope}[shift={(0,0,0)}]
    \IsoBlock{\LayerWidthA}{\LayerDepthA}{\LayerThicknessA}{\ColorLayerA}{}
    \node[iso label, anchor=east] at (-0.20,0.10,0.16) {bottom electrode};
    \draw[iso callout] (-0.10,0.12,0.16) -- (0.65,0.10,0.16);
  \end{scope}

  % Layer B: dielectric
  \begin{scope}[shift={(0,0,0.42)}]
    \IsoBlock{\LayerWidthA}{\LayerDepthA}{\LayerThicknessBD}{\ColorLayerBD}{}
    \node[iso label, anchor=east] at (-0.20,0.14,0.64) {dielectric polymer};
    \draw[iso callout] (-0.10,0.16,0.64) -- (0.80,0.15,0.64);
  \end{scope}

  % Layer C: charge region with trapped dots
  \begin{scope}[shift={(0,0,1.04)}]
    \IsoBlock{\LayerWidthA}{\LayerDepthA}{\LayerThicknessBCD}{\ColorLayerBCD}{}
    % Trapped charge dots: user can edit positions below
    \foreach \x/\y in {0.80/0.45,1.50/1.10,2.25/0.72,3.05/1.28,3.80/0.55}{
      \fill[iso charge] (\x,\y,0.20) circle[radius=0.055];
    }
    \node[iso label, anchor=west] at (5.95,-0.70,0.28) {trapped-charge region};
    \draw[iso callout] (5.75,-0.64,0.28) -- (3.78,0.58,0.20);
  \end{scope}

  % Layer D: top electrode
  \begin{scope}[shift={(0,0,1.56)}]
    \IsoBlock{\LayerWidthA}{\LayerDepthA}{\LayerThicknessBCDE}{\ColorLayerBCDE}{}
    \node[iso label, anchor=west] at (5.95,-0.88,0.32) {top electrode};
    \draw[iso callout] (5.75,-0.82,0.32) -- (4.25,0.35,0.12);
  \end{scope}

  % Depth cue: hidden edge
  \draw[iso hidden edge] (4.95,1.90,0.28) -- (4.95,1.90,1.78);
  \node[iso label, text=cGray!70, anchor=west] at (5.95,-0.50,0.95) {fixed-view depth cue};
\end{tikzpicture}

%% Or copy-paste the TikZ code below into your .tex and edit layer counts, colors, labels, charge positions.

\providecommand{\LayerCountA}{4}  % total layer count
\providecommand{\LayerThicknessA}{0.28}  % bottom electrode thickness (cm)
\providecommand{\LayerThicknessBD}{0.44}  % middle (dielectric) thickness (cm)
\providecommand{\LayerThicknessBCD}{0.34}  % charge layer thickness (cm)
\providecommand{\LayerThicknessBCDE}{0.22}  % top electrode thickness (cm)
\providecommand{\LayerWidthA}{4.8}  % all layers width (cm)
\providecommand{\LayerDepthA}{1.7}  % all layers depth (cm)
\providecommand{\ColorLayerA}{cAmber}  % bottom electrode
\providecommand{\ColorLayerBD}{cTeal}  % dielectric polymer
\providecommand{\ColorLayerBCD}{cBlue}  % charge layer
\providecommand{\ColorLayerBCDE}{cGray}  % top electrode

\begin{tikzpicture}[iso basis]
  \node[font=\fontsize{9}{10.8}\selectfont\bfseries, text=cGray!90!black, anchor=west]
    at (-0.40, -1.10, 2.95) {Layered device stack};

  % Layer A: bottom electrode
  \begin{scope}[shift={(0,0,0)}]
    \IsoBlock{\LayerWidthA}{\LayerDepthA}{\LayerThicknessA}{\ColorLayerA}{}
    \node[iso label, anchor=east] at (-0.20,0.10,0.16) {bottom electrode};
    \draw[iso callout] (-0.10,0.12,0.16) -- (0.65,0.10,0.16);
  \end{scope}

  % Layer B: dielectric
  \begin{scope}[shift={(0,0,0.42)}]
    \IsoBlock{\LayerWidthA}{\LayerDepthA}{\LayerThicknessBD}{\ColorLayerBD}{}
    \node[iso label, anchor=east] at (-0.20,0.14,0.64) {dielectric polymer};
    \draw[iso callout] (-0.10,0.16,0.64) -- (0.80,0.15,0.64);
  \end{scope}

  % Layer C: charge region with trapped dots
  \begin{scope}[shift={(0,0,1.04)}]
    \IsoBlock{\LayerWidthA}{\LayerDepthA}{\LayerThicknessBCD}{\ColorLayerBCD}{}
    % Trapped charge dots: user can edit positions below
    \foreach \x/\y in {0.80/0.45,1.50/1.10,2.25/0.72,3.05/1.28,3.80/0.55}{
      \fill[iso charge] (\x,\y,0.20) circle[radius=0.055];
    }
    \node[iso label, anchor=west] at (5.95,-0.70,0.28) {trapped-charge region};
    \draw[iso callout] (5.75,-0.64,0.28) -- (3.78,0.58,0.20);
  \end{scope}

  % Layer D: top electrode
  \begin{scope}[shift={(0,0,1.56)}]
    \IsoBlock{\LayerWidthA}{\LayerDepthA}{\LayerThicknessBCDE}{\ColorLayerBCDE}{}
    \node[iso label, anchor=west] at (5.95,-0.88,0.32) {top electrode};
    \draw[iso callout] (5.75,-0.82,0.32) -- (4.25,0.35,0.12);
  \end{scope}

  % Depth cue: hidden edge
  \draw[iso hidden edge] (4.95,1.90,0.28) -- (4.95,1.90,1.78);
  \node[iso label, text=cGray!70, anchor=west] at (5.95,-0.50,0.95) {fixed-view depth cue};
\end{tikzpicture}

%% Or copy-paste the TikZ code below into your .tex and edit layer counts, colors, labels, charge positions.

\providecommand{\LayerCountA}{4}  % total layer count
\providecommand{\LayerThicknessA}{0.28}  % bottom electrode thickness (cm)
\providecommand{\LayerThicknessBD}{0.44}  % middle (dielectric) thickness (cm)
\providecommand{\LayerThicknessBCD}{0.34}  % charge layer thickness (cm)
\providecommand{\LayerThicknessBCDE}{0.22}  % top electrode thickness (cm)
\providecommand{\LayerWidthA}{4.8}  % all layers width (cm)
\providecommand{\LayerDepthA}{1.7}  % all layers depth (cm)
\providecommand{\ColorLayerA}{cAmber}  % bottom electrode
\providecommand{\ColorLayerBD}{cTeal}  % dielectric polymer
\providecommand{\ColorLayerBCD}{cBlue}  % charge layer
\providecommand{\ColorLayerBCDE}{cGray}  % top electrode

\begin{tikzpicture}[iso basis]
  \node[font=\fontsize{9}{10.8}\selectfont\bfseries, text=cGray!90!black, anchor=west]
    at (-0.40, -1.10, 2.95) {Layered device stack};

  % Layer A: bottom electrode
  \begin{scope}[shift={(0,0,0)}]
    \IsoBlock{\LayerWidthA}{\LayerDepthA}{\LayerThicknessA}{\ColorLayerA}{}
    \node[iso label, anchor=east] at (-0.20,0.10,0.16) {bottom electrode};
    \draw[iso callout] (-0.10,0.12,0.16) -- (0.65,0.10,0.16);
  \end{scope}

  % Layer B: dielectric
  \begin{scope}[shift={(0,0,0.42)}]
    \IsoBlock{\LayerWidthA}{\LayerDepthA}{\LayerThicknessBD}{\ColorLayerBD}{}
    \node[iso label, anchor=east] at (-0.20,0.14,0.64) {dielectric polymer};
    \draw[iso callout] (-0.10,0.16,0.64) -- (0.80,0.15,0.64);
  \end{scope}

  % Layer C: charge region with trapped dots
  \begin{scope}[shift={(0,0,1.04)}]
    \IsoBlock{\LayerWidthA}{\LayerDepthA}{\LayerThicknessBCD}{\ColorLayerBCD}{}
    % Trapped charge dots: user can edit positions below
    \foreach \x/\y in {0.80/0.45,1.50/1.10,2.25/0.72,3.05/1.28,3.80/0.55}{
      \fill[iso charge] (\x,\y,0.20) circle[radius=0.055];
    }
    \node[iso label, anchor=west] at (5.95,-0.70,0.28) {trapped-charge region};
    \draw[iso callout] (5.75,-0.64,0.28) -- (3.78,0.58,0.20);
  \end{scope}

  % Layer D: top electrode
  \begin{scope}[shift={(0,0,1.56)}]
    \IsoBlock{\LayerWidthA}{\LayerDepthA}{\LayerThicknessBCDE}{\ColorLayerBCDE}{}
    \node[iso label, anchor=west] at (5.95,-0.88,0.32) {top electrode};
    \draw[iso callout] (5.75,-0.82,0.32) -- (4.25,0.35,0.12);
  \end{scope}

  % Depth cue: hidden edge
  \draw[iso hidden edge] (4.95,1.90,0.28) -- (4.95,1.90,1.78);
  \node[iso label, text=cGray!70, anchor=west] at (5.95,-0.50,0.95) {fixed-view depth cue};
\end{tikzpicture}

%% Or copy-paste the TikZ code below into your .tex and edit layer counts, colors, labels, charge positions.

\providecommand{\LayerCountA}{4}  % total layer count
\providecommand{\LayerThicknessA}{0.28}  % bottom electrode thickness (cm)
\providecommand{\LayerThicknessBD}{0.44}  % middle (dielectric) thickness (cm)
\providecommand{\LayerThicknessBCD}{0.34}  % charge layer thickness (cm)
\providecommand{\LayerThicknessBCDE}{0.22}  % top electrode thickness (cm)
\providecommand{\LayerWidthA}{4.8}  % all layers width (cm)
\providecommand{\LayerDepthA}{1.7}  % all layers depth (cm)
\providecommand{\ColorLayerA}{cAmber}  % bottom electrode
\providecommand{\ColorLayerBD}{cTeal}  % dielectric polymer
\providecommand{\ColorLayerBCD}{cBlue}  % charge layer
\providecommand{\ColorLayerBCDE}{cGray}  % top electrode

\begin{tikzpicture}[iso basis]
  \node[font=\fontsize{9}{10.8}\selectfont\bfseries, text=cGray!90!black, anchor=west]
    at (-0.40, -1.10, 2.95) {Layered device stack};

  % Layer A: bottom electrode
  \begin{scope}[shift={(0,0,0)}]
    \IsoBlock{\LayerWidthA}{\LayerDepthA}{\LayerThicknessA}{\ColorLayerA}{}
    \node[iso label, anchor=east] at (-0.20,0.10,0.16) {bottom electrode};
    \draw[iso callout] (-0.10,0.12,0.16) -- (0.65,0.10,0.16);
  \end{scope}

  % Layer B: dielectric
  \begin{scope}[shift={(0,0,0.42)}]
    \IsoBlock{\LayerWidthA}{\LayerDepthA}{\LayerThicknessBD}{\ColorLayerBD}{}
    \node[iso label, anchor=east] at (-0.20,0.14,0.64) {dielectric polymer};
    \draw[iso callout] (-0.10,0.16,0.64) -- (0.80,0.15,0.64);
  \end{scope}

  % Layer C: charge region with trapped dots
  \begin{scope}[shift={(0,0,1.04)}]
    \IsoBlock{\LayerWidthA}{\LayerDepthA}{\LayerThicknessBCD}{\ColorLayerBCD}{}
    % Trapped charge dots: user can edit positions below
    \foreach \x/\y in {0.80/0.45,1.50/1.10,2.25/0.72,3.05/1.28,3.80/0.55}{
      \fill[iso charge] (\x,\y,0.20) circle[radius=0.055];
    }
    \node[iso label, anchor=west] at (5.95,-0.70,0.28) {trapped-charge region};
    \draw[iso callout] (5.75,-0.64,0.28) -- (3.78,0.58,0.20);
  \end{scope}

  % Layer D: top electrode
  \begin{scope}[shift={(0,0,1.56)}]
    \IsoBlock{\LayerWidthA}{\LayerDepthA}{\LayerThicknessBCDE}{\ColorLayerBCDE}{}
    \node[iso label, anchor=west] at (5.95,-0.88,0.32) {top electrode};
    \draw[iso callout] (5.75,-0.82,0.32) -- (4.25,0.35,0.12);
  \end{scope}

  % Depth cue: hidden edge
  \draw[iso hidden edge] (4.95,1.90,0.28) -- (4.95,1.90,1.78);
  \node[iso label, text=cGray!70, anchor=west] at (5.95,-0.50,0.95) {fixed-view depth cue};
\end{tikzpicture}
